\subsection*{NT-E14--E16. Finite-field detection of characteristic-zero isogeny}

The analytic input here is Dedekind RH for certain finite
number fields. It is not RH for the standard $L$-function
of a non-CM elliptic curve.
For an even squarefree integer $m\ge2$, let
$\mathcal D_m$ be the collection of finite Galois
extensions $T/\mathbf Q$ satisfying
\[
 [T:\mathbf Q]\le384,\qquad
 T/\mathbf Q\text{ is unramified at every prime }p\nmid m.
 \tag{E14.1}
\]
The hypothesis RH for $\zeta_{\mathbf Q}$ and for
$\zeta_T$, $T\in\mathcal D_m$, is sufficient for the
applications below. This is a finite family: the
discriminant bound (E16.4) below and Hermite's theorem
give finiteness. A precise reference is Milne,
\emph{Algebraic Number Theory}, v3.08 (2020),
Theorems~8.42--8.43, printed pp.~151--152.
Input $M$ supplies this family.
The more restricted fields actually used in the
argument are specified explicitly below.

\paragraph{NT-E14: the published non-CM isogeny test.}
Let $E_1,E_2/\mathbf Q$ be elliptic curves without
complex multiplication, with conductors $N_1,N_2$,
and put $m=\operatorname{rad}(2N_1N_2)$.
If the two curves are not $\mathbf Q$-isogenous,
then there is a prime $p\nmid2N_1N_2$ with
\[
 a_p(E_1)\ne a_p(E_2),\qquad
 p\le \bigl(482\log m+2880\bigr)^2.
 \tag{E14.2}
\]
Consequently equality of the point counts at all
such primes up to this bound is equivalent to
$\mathbf Q$-isogeny.
This is Mayle--Wang, \emph{On the effective version
of Serre's open image theorem},
Bull.\ London Math.\ Soc.~56 (2024), 1399--1416,
Theorem~2, printed p.~1401; complete proof
pp.~1408--1410. Its proof omits $2$ as well as all
bad primes, although the theorem statement only
requires common good reduction.

The precise source dependency is useful. Proposition~12,
pp.~1407--1408, constructs a finite Galois quotient field
\[
 F\subset\mathbf Q(E_1[2^\infty],E_2[2^\infty]),
                      \qquad [F:\mathbf Q]\le192.
 \tag{E14.3}
\]
It detects a trace difference. Corollary~6,
pp.~1404--1405, then applies effective Chebotarev to
$T=F(\sqrt m)$ to exclude every prime dividing $m$.
Thus only RH for this actual $\zeta_T$ and
$\zeta_{\mathbf Q}$ is used analytically.
The field $F$ is a quotient of the full $2$-adic image;
it is not asserted to be the first $2$-division field.
The convenient fixed family (E14.1) contains every
possible such $T$. The improved numerical constants
in (E14.2) use the published Bach--Sorenson Table~3
bounds through the source's computation
(13)--(16). Those numerical inputs are credited,
not represented as a fresh rerun of that table.

\paragraph{NT-E15: the published effective open-image bound.}
For a non-CM elliptic curve $E/\mathbf Q$ with
conductor $N_E$, input $M$ implies that
$\overline\rho_{E,\ell}$ is surjective for every
prime $\ell>B_E$, where
\[
                 B_E=964\log\operatorname{rad}(2N_E)+5760.
 \tag{E15.1}
\]
This is Mayle--Wang Theorem~1, printed p.~1401,
proved in \S5.2, pp.~1412--1413.
It retains both the base field $\mathbf Q$ and
the non-CM hypothesis. It does not prove the
uniform proposed cutoff $37$ for all curves.

Here a sufficient analytic family is (E14.1) for
$m=\operatorname{rad}(2N_E)$, together with
RH for $\zeta_{\mathbf Q}$. Indeed the proof
associates to a nonsurjective prime $\ell>37$ a
nontrivial quadratic twist $E_D$ with
$\operatorname{rad}(N_{E_D})\mid m$.
It applies (E14.2) to $E,E_D$; the two curves
are not $\mathbf Q$-isogenous by its Lemma~14.
The witnessing prime has a nonzero trace divisible
by $\ell$, and the Hasse bound gives
$\ell\le2\sqrt p$, producing (E15.1).
All analytic fields are therefore the fields
$F_D(\sqrt m)$ from (E14.3) for these twists,
of degree at most $384$ and unramified outside $m$.
The finitely many relevant squarefree $D$ divide
$N_E$, with their signs allowed.
The source's modular-curve exclusions, local
inertia lemma, Faltings theorem, and twist facts
are unconditional classical inputs; this reduction
does not claim to reprove them.

\paragraph{NT-E16: a complete decision formulation from integral equations.}
Let $E_1,E_2$ be given by nonsingular integral
Weierstrass equations over $\mathbf Q$, of
discriminants $\Delta_1,\Delta_2\ne0$.
No minimality, factorization, conductor, or
non-CM promise is required in this paragraph.
Put
\[
 R=2|\Delta_1\Delta_2|,\qquad
 b=\operatorname{bitlength}(R),\qquad
                      B=(1532b+10181)^2.
 \tag{E16.1}
\]
Under $M$, the curves are $\mathbf Q$-isogenous
if and only if
\[
 \#E_1(\mathbf F_p)=\#E_2(\mathbf F_p)
                  \quad\hbox{for every prime }p\le B,\ p\nmid R.
 \tag{E16.2}
\]
Testing (E16.2) by direct point counting is a
deterministic algorithm polynomial in the total
bit length of the integral coefficients.
This decides isogeny over $\mathbf Q$; it does
not construct an isogeny, decide geometric
isogeny over arbitrary extensions, or compute an
endomorphism ring.

More precisely, the sufficient RH family is
RH for $\zeta_{\mathbf Q}$ and
all $\zeta_T$ in (E14.1) with
$m=\operatorname{rad}(R)$.
The algorithm never needs to compute $m$.
For any nonisogenous pair only one actual
quotient field $F$ and the field $F(\sqrt m)$
are needed in the proof.

\begin{proof}
We give the finite-quotient argument explicitly;
it is the $\ell=2$, rank-two case of the
Faltings--Serre construction in Mayle--Wang
Proposition~12.
Faltings's isogeny theorem says that two abelian
varieties are isogenous over their number field
if and only if their local Euler factors agree
at all but finitely many places: Invent.\ Math.~73
(1983), Corollary~2, equivalence (i)--(iii),
printed p.~361. For elliptic curves over
$\mathbf Q$, the good local factor is determined
by $1-a_pT+pT^2$. If $E_1,E_2$ are not
$\mathbf Q$-isogenous, we can therefore choose
a common good odd prime where their traces differ.
A Frobenius element at that prime has differing
rational $2$-adic Tate traces. This argument also
applies to CM curves and needs no separate
character-criterion citation. Mayle--Wang recall
the elliptic criterion in (11), printed p.~1407.

Write $\rho=(\rho_1,\rho_2)$ for the product
of the integral rank-two Tate representations.
Let $A$ be the $\mathbf Z_2$-linear span of
$\rho(G_{\mathbf Q})$ in
$M_2(\mathbf Z_2)\times M_2(\mathbf Z_2)$.
Products of image elements remain image elements,
so $A$ is a ring containing the identity.
It is a finite free $\mathbf Z_2$-module of rank
at most eight. Let $e$ be the least finite
$2$-adic valuation of a nonzero difference
$\operatorname{tr}\rho_1(\sigma)-
 \operatorname{tr}\rho_2(\sigma)$.
The linear functional
\[
 \lambda:A\longrightarrow\mathbf Z_2,\qquad
 (u,v)\longmapsto
             2^{-e}(\operatorname{tr}u-\operatorname{tr}v)
 \tag{E16.3}
\]
is integral on $A$ and induces a nonzero functional
on at least one image element in $A/2A$.
Let $G$ be the image of $G_{\mathbf Q}$ in
$(A/2A)^\times$, and let $C$ be the nonempty
subset on which this functional is nonzero.
Trace invariance makes $C$ conjugacy-stable.
Every lift of an element of $C$ has different
two Tate traces.

Here $|G|\le2^8-1=255$: zero is not a unit.
The homomorphism to $G$ is continuous.
Indeed $A$ is a closed finite lattice in the
ambient matrix module, and $2A$ is open in $A$.
The fixed field $F$ of its kernel is consequently
a finite Galois extension of $\mathbf Q$,
contained in the joint $2$-adic division tower.
It lies in some finite level
$\mathbf Q(E_1[2^j],E_2[2^j])$, with $j\ge1$, so its degree
divides
$|\mathrm{GL}_2(\mathbf Z/2^j\mathbf Z)|^2
 =36\cdot2^{8j-8}$.
Among divisors of such an integer not exceeding
$255$, the largest is $192$.
Thus $[F:\mathbf Q]\le192$.
This argument uses $A/2A$ itself; it does not
identify it with the reduction of the ambient
matrix algebra, an identification that need
not be injective for an unsaturated lattice.

N\'eron--Ogg--Shafarevich gives that $F/\mathbf Q$
is unramified outside the prime divisors of $R$.
Here the exact cited form is Mayle--Wang
Theorem~8 and Corollary~9, printed pp.~1406--1407,
with their Remark~11 for a product of two curves;
Theorem~8 cites Serre--Tate, \emph{Good reduction
of abelian varieties}, Ann.\ Math.~88 (1968),
Theorem~1.
Set $m=\operatorname{rad}(R)$ and
$T=F(\sqrt m)$. This is Galois of degree
$d\le384$ and unramified outside $m$.
As $m$ is even and squarefree, every prime
dividing $m$ ramifies in $\mathbf Q(\sqrt m)$,
and hence in $T/\mathbf Q$.
The inverse image of $C$ in
$\operatorname{Gal}(T/\mathbf Q)$ is nonempty
and conjugacy-stable.
NT-C03 applied to any one of its conjugacy
classes provides, under RH for $\zeta_T$
and $\zeta_{\mathbf Q}$, an unramified prime
$p$ with differing Tate traces and
\[
               p\le(4\log D_T+2.5d+5)^2.
\]
In particular $p\nmid R$, so both input equations
have good reduction and the differing traces
are their actual $a_p(E_i)$.

The standard different bound, in the exact
form of Mayle--Wang Lemma~7, printed p.~1405,
gives
\[
 \begin{split}
 \log D_T
   &\le(d-1)\log\operatorname{rad}(D_T)+d\log d\\
   &\le383\log m+384\log384.
 \end{split}
 \tag{E16.4}
\]
For completeness the local calculation is
$v_{\mathfrak q}(\mathfrak D_T)
 \le e_{\mathfrak q}-1+
             e_{\mathfrak q}v_p(e_{\mathfrak q})$.
For a Galois field, $e_{\mathfrak q}\mid d$;
summing residue degrees over primes above $p$
bounds the exponent of $p$ in $D_T$ by
$d-1+d\,v_p(d)$.
Sum after multiplying by $\log p$ over
ramified $p$. The first terms contribute
$(d-1)\log\operatorname{rad}(D_T)$ and the
remaining terms are at most $d\log d$.
This proves (E16.4), retaining the classical
local different inequality as its input.

Since $\log384<6$, the Chebotarev expression is
bounded by
\[
 \begin{split}
 4\log D_T+2.5d+5
    &\le1532\log m+1536\log384+965\\
    &<1532\log R+10181
     \le1532b+10181.
 \end{split}
 \tag{E16.5}
\]
The last inequality uses $R<2^b$ and $\log2<1$.
This proves the existence of a differing point
count within (E16.1) whenever the two curves
are not $\mathbf Q$-isogenous.
Conversely a $\mathbf Q$-isogeny identifies
the rational Tate representations, so the
point counts agree at every common good
prime. This proves the equivalence.

Finally let $L$ be the total input bit length.
Each discriminant is a fixed polynomial in
the five integral Weierstrass coefficients;
it is computable in polynomial time and
has bit length $O(L)$. Thus $b=O(L)$ and
$B=O(L^2)$. Trial division or a sieve
enumerates all primes up to the integer $B$
in polynomial time. Testing $p\nmid R$ only
needs a remainder calculation.
At each retained prime enumerate all
$(x,y)\in\mathbf F_p^2$ satisfying the reduced
Weierstrass equation and add the unique point
at infinity. This uses $O(p^2)$ elementary
modular operations; even the sum over all
integers $p\le B$ is $O(B^3)$.
The operands have polynomial bit length.
Consequently the entire finite test is
polynomial in $L$, with no factoring of $R$,
conductor calculation, or point-counting
oracle. A negative output includes a prime
and two different counts as an unconditional
certificate; correctness of every positive
output uses the stated RH input.
\end{proof}

NT-E14--E15 are precise published theorem
transfers with their analytic field families
identified from the full source proofs.
NT-E16 is a fully proved formulation of the
finite-quotient method with a weaker elementary
constant and a direct input model; no
literature-priority claim is made.
Geometrically, this detects isomorphism of
the semisimple rank-two Tate local systems
using closed points of the good-reduction
open of $\operatorname{Spec}\mathbf Z$,
then invokes Faltings to recover
$\mathbf Q$-isogeny. The finite quotients
have nonabelian images in general, but their
Dedekind functions are trivial-character
instances of $M$. No non-CM $\mathrm{GL}_2$
GRH or general Artin holomorphy has entered.
Work/run: \texttt{effective\_elliptic\_detection.tex},
\texttt{GRH\_NIGHT/number\_theory}.
